% ============================================================================
% CHAPTER 1 — INTRODUCTION
% ============================================================================

\chapter{Introduction}
\label{ch:intro}

\section{The Problem}

Supply chains run on four enterprise systems: \gls{srm} (procurement), \gls{wms} (inventory), \gls{oms} (sales), \gls{tms} (shipping). These systems come from different vendors and do not talk to each other. A GEODIS survey reports only 6\% of companies achieve full visibility across them~\cite{geodis2024}; operations managers spend 15 to 30 minutes per investigation manually stitching data from three or four interfaces~\cite{aws2024sync,omniful2026}.

Most cross-system questions are blocked by this fragmentation: \emph{why is this order late?}, \emph{which customers are at risk if product X runs out?}, \emph{is the chosen carrier appropriate for that product?}

\section{The Question}

\emph{Can an AI agent with read-only access to the four systems answer realistic operational questions at human-level factual accuracy, while significantly reducing investigation time?}

\section{SERGE}

We build \textbf{SERGE}, a read-only agent that sits over the four supply chain systems via the Model Context Protocol~\cite{mcp2024}. A simulated B2B mango distributor is used as the target domain; its four systems are exposed as 35 read-only \gls{mcp} tools. Claude Code acts as the reasoning agent, invoked from the user's desktop client.

The system is evaluated on 27 scenarios spanning five complexity levels (L1 trivial to L5 expert). Each scenario is run three times, judged by a separate \gls{llm} on four dimensions (factual accuracy, completeness, hallucination, reasoning), and aggregated. Two passes are reported: a baseline and a refined run after one iteration cycle.

\section{Contributions}

\begin{enumerate}
    \item A plug-and-play \gls{mcp} architecture for supply chain systems, agent-agnostic by design.
    \item A read-only diagnostic agent pattern, complementary to the optimisation-centric supply chain \gls{ai} literature.
    \item An \gls{llm}-as-judge methodology with an explicit \emph{flexible-on-form, strict-on-facts} principle and a structured failure taxonomy.
    \item A reusable benchmark of 27 scenarios at five complexity levels.
\end{enumerate}

\section{Headline Results}

Refined pass (Run~2): 94\% pass rate over 81 evaluations, weighted score 4.84/5.0, \textbf{zero hallucinations}, 100\% pass on L4 (complex cross-system) scenarios. Resolution time averages about 13\,s on L1 and 72\,s on L5.

\section{Report Map}

Chapter~\ref{ch:background}: background on \gls{llm}s, agents, and \gls{mcp}. Chapter~\ref{ch:relwork}: related work. Chapter~\ref{ch:modelisation}: the modelled business, the entity model, and how the dataset was seeded. Chapter~\ref{ch:architecture}: the agent--\gls{mcp}--backend stack. Chapter~\ref{ch:methodology}: evaluation methodology. Chapter~\ref{ch:scenarios}: the 27 scenarios. Chapter~\ref{ch:results}: empirical results, including the refinement cycle. Chapter~\ref{ch:discussion}: discussion. Chapter~\ref{ch:conclusion}: conclusion.
